% فرم پروپوزال دانشگاه تربیت مدرس - دانشکده علوم ریاضی - علوم کامپیوتر
% مقطع کارشناسی ارشد

% راهنما
 
% بخش های مختلف، به صورت فایل های جداگانه آورده شده اند.
% این بخش ها در پوشه 
% Segments
% قرار دارند.

% به طور مثال جدول ۱ که شامل اطلاعات دانشجو است به صورت فایل 
% Segment-1
% آورده شده است.

% برای پر کردن جداول، برای هر بخش مشخص، به فایل 
% Input- 
% آن مراجعه کنید.
% این فایل ها در پوشه 
% Inputs
% قرار دارند.

% به طور مثال برای پر کردن جدول ۱ ، فایل 
% (Segment-1)
% کافیست فایل 
% (Segment-1-Input)
% را باز کرده و با توجه به محتوای آن، خانه ها را پر کنید.
    
% به طور خلاصه، برای پر کردن جداول و محتوای پروپوزال، کافیست فایل های موجود در پوشه 
% Inputs
% را باز کرده و پر کنید.

% راهنما 


\documentclass[a4paper, 12pt]{article}

\input{packages}
\input{commands}

\input{Inputs/Segment-1-Input}
\input{Inputs/Segment-2-Input}
\input{Inputs/Segment-3-Input}

\input{Inputs/Introduction-Input}
\input{Inputs/Theories-Input}
\input{Inputs/Method-Input}
\input{Inputs/Innovation-Input}
\input{Inputs/Reference-Input}

\input{Inputs/Segment-4-Input}
\input{Inputs/Segment-5-Input}
\input{Inputs/Segment-6-Input}
\input{Inputs/Segment-7-Input}
\input{Inputs/Segment-8-Input}


\begin{document}
	
	\input{Header}  
	
	\input{Segments/Segment-1} % شامل جدول اول - اطلاعات دانشجو
	\input{Segments/Segment-2} % شامل جدول دوم - اطلاعات اساتید راهنما و مشاور
	\input{Segments/Segment-3} % شامل جدول سوم - موضوع پروپوزال
	
	\input{Inputs/Sections-Input} % شامل بخش های مختلف پروپوزال
	
	\newpage
	\input{Segments/Segment-4} % شامل جدول چهار - موضوع تصویب شده
	\input{Segments/Segment-5} % شامل جدول پنج - اطلاعات کمیته گروه
	\input{Segments/Segment-6} % شامل جدول شش - توضیحات
	\input{Segments/Segment-7} % شامل جدول هفت - امضا مدیر گروه
	\input{Segments/Segment-8} % شامل جدول هشت - امضا رئیش شورا
	
	
\end{document}
